\chapter{Myelofibrosis}
\index{Myelofibrosis}
\label{sec:Myelofibrosis}

\section{Overview}
Myelofibrosis is a myeloproliferative tumor characterized by clonal megakaryocyte growth, bone marrow scarring (reticulin and/or collagen deposition), extramedullary hematopoiesis (splenomegaly, hepatomegaly), and progressive cytopenias, which may arise de novo (primary MF) or evolve from PV or ET (post-PV/post-ET MF). The condition is driven by JAK2, CALR, or MPL mutations in over 90\% of cases, with ASXL1, EZH2, SRSF2 mutations associated with worse outlook. This degenerative condition involves progressive deterioration of affected tissues, leading to gradual loss of function. It is more common with advancing age.
\section{Causes}
\section{Risk Factors}

\begin{itemize}[noitemsep]
  \item Genetic predisposition and family history
  \item Advancing age
  \item Lifestyle factors (diet, exercise, smoking, alcohol)
  \item Environmental exposures
  \item Underlying medical conditions
  \item Medication use
\end{itemize}
\section{Signs and Symptoms}

\subsection{Common symptoms}
\begin{itemize}[noitemsep]
  \item \item You may experience massive splenomegaly (early satiety, abdominal discomfort), constitutional symptoms (fatigue, weight loss, night sweats, fever), bone pain, and progressive anemia requiring transfusions. Pain or discomfort in the affected area    \item Pain or discomfort in the affected area
  \end{itemize}

\subsection{Emergency warning signs}
\begin{itemize}[noitemsep]
  \item Severe or worsening symptoms of myelofibrosis require immediate medical evaluation
\end{itemize}
\section{Progression and Stages}
The condition may progress through different stages of severity over time. Early stages often involve mild or subtle symptoms that may be overlooked. Moderate stages involve more noticeable symptoms affecting daily function. Advanced stages may involve significant impairment and complications.
\section{Complications}
\begin{itemize}[noitemsep]
  \item Disease progression leading to worsening symptoms
  \item Secondary infections or injuries
  \item Side effects of treatment
  \item Reduced quality of life
  \item Emotional and psychological impact
\end{itemize}
\section{Diagnosis}
To make the diagnosis, doctors look for bone marrow biopsy showing megakaryocyte growth and atypia, reticulin/collagen scarring, and exclusion of other causes. The Dynamic International Prognostic Scoring System (DIPSS) and molecular-enhanced versions (MIPSS70) guide treatment. Management options include ruxolitinib (JAK1/2 inhibitor, improves splenomegaly and symptoms), fedratinib (JAK2 inhibitor), momelotinib (JAK1/2 inhibitor with anemia benefit), and allogeneic stem cell transplantation (the only curative option). Comprehensive medical history and physical examination Laboratory tests to assess organ function and rule out other conditions Imaging studies to visualize affected structures Specialized testing depending on the affected system Consultation with relevant specialists Monitoring of disease progression over time

\begin{itemize}[noitemsep]
  \item Comprehensive medical history and physical examination
  \item Laboratory tests to assess organ function
  \item Imaging studies to visualize affected structures
  \item Specialized testing depending on the affected system
  \item Consultation with relevant specialists
\end{itemize}
\section{Prevention}
\begin{itemize}[noitemsep]
  \item Maintain a healthy lifestyle with balanced diet and exercise
  \item Avoid known risk factors
  \item Attend regular medical check-ups for early detection
  \item Follow recommended screening guidelines
\end{itemize}
\section{Treatment Options}
Medications to manage symptoms and address underlying mechanisms Lifestyle modifications and self-care measures Physical therapy and rehabilitation Surgical intervention when indicated Regular monitoring and follow-up care Patient education and support services

\begin{itemize}[noitemsep]
  \item Medications to manage symptoms and address underlying mechanisms
  \item Lifestyle modifications and self-care measures
  \item Physical therapy and rehabilitation
  \item Surgical intervention when indicated
  \item Regular monitoring and follow-up care
\end{itemize}
\section{Living with It}
\begin{itemize}[noitemsep]
  \item Follow your treatment plan as prescribed
  \item Maintain regular follow-up with your healthcare provider
  \item Make necessary lifestyle adjustments
  \item Seek support from family, friends, or support groups
  \item Stay informed about your condition
\end{itemize}
\section{Outlook}
The outlook varies from person to person depending on risk category. The outlook varies depending on the specific condition, severity at diagnosis, and response to treatment. Early intervention generally leads to better outcomes. Many conditions can be effectively managed with appropriate treatment. Regular follow-up care is important for monitoring and treatment adjustment.
